\section{Methodology}
\label{sec:methodology}

\subsection{Dataset and Preprocessing}
\label{sec:method-data}

We use the same four-class dataset as Tasso et al.~\cite{tasso2026lesion}: 864
ultra-widefield fundus images labelled healthy, choroidal nevus, uveal melanoma,
or CHRPE, with 302, 290, 170, and 102 images per class. Uveal melanoma and CHRPE
are the least represented, and CHRPE, the smallest class, is also the hardest to
classify, as it overlaps visually with nevus. The test set holds 133 images, with
47 healthy, 44 nevus, 26 melanoma, and 16 CHRPE.

Their pipeline resizes each image to $512\times512$ and normalizes it with
ImageNet statistics, and the teacher and assistant networks
(Section~\ref{sec:method-kd}) keep this unchanged. The deployed student instead
runs at $160\times160$, the native resolution of its backbone, reached by trimming
the dark border around the circular field of view and resizing; the normalization
is the same.

\begin{figure*}[t]
\centering
\resizebox{0.95\textwidth}{!}{%
\begin{tikzpicture}[
    font=\footnotesize,
    stage/.style={
        draw=black!75,
        rounded corners=2pt,
        align=center,
        minimum width=2.55cm,
        minimum height=0.78cm,
        inner xsep=5pt,
        inner ysep=4pt
    },
    flow/.style={-{Latex[length=2.2mm,width=1.5mm]}, line width=0.35pt},
    qattrain/.style={flow, rounded corners=5pt},
    init/.style={flow, dashed},
    edge note/.style={font=\scriptsize, inner sep=0pt}
]
\node[stage] (crop) {ResNet-50\\lesion crops};
\node[stage, right=0.65cm of crop] (field) {ResNet-50\\full field};
\node[stage, right=0.65cm of field] (assistant) {ResNet-18\\assistant};
\node[stage, right=0.65cm of assistant] (student) {\texttt{test\_resnet}\\FP32 student};
\node[stage, right=0.65cm of student] (quantized) {quantized\\\texttt{test\_resnet}};

\draw[flow] (crop) -- node[midway, above=4pt, edge note] {KD} (field);
\draw[flow] (field) -- node[midway, above=4pt, edge note] {KD} (assistant);
\draw[flow] (assistant) -- node[midway, above=4pt, edge note] {KD} (student);
\draw[init] (student) -- node[midway, above=4pt, edge note] {init.} (quantized);
\draw[qattrain] (assistant.south) -- ++(0,-0.85)
    -- node[midway, above=4pt, edge note] {KD}
    ($(quantized.south)+(0,-0.85)$) -- (quantized.south);
\end{tikzpicture}%
}
\caption{The teacher-assistant distillation chain.}
\label{fig:takd}
\end{figure*}

\subsection{Teacher-Assistant Knowledge Distillation}
\label{sec:method-kd}

The deployed network must be small enough to fit the FPGA, which constrains the
architecture before accuracy is considered. On a dataset this small, training from
scratch forces a bad trade: a large network overfits, and a small one learns too
little. ImageNet-pretrained initialization avoids it by giving a small network
visual features it could not acquire from the data alone. For the deployable
student we need this initialization in a backbone much smaller than the ResNet
teachers. We therefore instantiate the student with timm using the pretrained
\texttt{test\_resnet.r160\_in1k} checkpoint hosted on Hugging Face
\cite{hfTestResnet}: a ResNet-style model with native $160\times160$ resolution
and ImageNet-1k weights. Fine-tuning it on the four-class task
improves on the small from-scratch baselines but plateaus short of the accuracy
we require, which is what motivates distillation.

Following Tasso et al., we first train a ResNet-50 \cite{he2016resnet} on
lesion-centered crops and use it to teach a ResNet-50 operating on the full
fundus field. This full-field ResNet-50 is the large teacher used for the
remaining chain. Distilling it directly into \texttt{test\_resnet} did not reach
the accuracy we expected, which is unsurprising when teacher and student are this
far apart in capacity: the student cannot match the much larger teacher's output
distribution \cite{cho2019efficacy}. The established remedy is to bridge the gap
with an intermediate network \cite{mirzadeh2020takd}, so we distil the full-field
ResNet-50 into a ResNet-18 and make that the teacher of \texttt{test\_resnet}.
Figure~\ref{fig:takd} shows the resulting teacher-assistant chain.

Every stage minimizes the same objective, a convex combination of the
cross-entropy on the labels and a Kullback-Leibler term that aligns the student's
softened output with the teacher's,
\begin{equation}
\mathcal{L} = \alpha\,\mathcal{L}_{\mathrm{CE}}(y, p_s)
            + (1-\alpha)\,\mathcal{L}_{\mathrm{KL}}(p_t^{\tau}, p_s^{\tau}),
\end{equation}
with $p_t^{\tau}$ and $p_s^{\tau}$ the teacher and student outputs softened at
temperature $\tau$, and $\alpha$ the weight on the two terms. We keep
$\alpha = 0.25$ and $\tau = 3$ throughout, including the quantization-aware
training of Section~\ref{sec:method-qat}, where distillation stays active.

\subsection{Quantization-Aware Training}
\label{sec:method-qat}

For deployment the fine-tuned student is reimplemented in Brevitas
\cite{pappalardo2023brevitas}, its convolutions, activations, and final linear
layer replaced by quantized counterparts that constrain weights and activations to
a fixed bit-width. We initialize this quantized model from the FP32 student and
run a calibration pass to set the quantizer scales. Loading the weights transfers
the trained floating-point solution into the Brevitas model; calibration then
defines how activations, and the already-loaded weights, are mapped to the
available fixed-point grid. We then continue training with quantizers active. The
distillation objective and the same teacher remain in use, so the network adapts
its weights to the precision it will use at inference instead of relying only on
the calibrated quantized warm start.

We train one model per precision, quantizing weights and activations to the same
width at 8, 6, and 4 bits. Lower precision yields a smaller accelerator but
degrades accuracy, a trade-off we examine in Section~\ref{sec:res-qat}.

\subsection{Hardware Compilation and Deployment}
\label{sec:method-hw}

The trained quantized model is exported from Brevitas to QONNX and handed to FINN,
which compiles it into the streaming dataflow accelerator described in
Section~\ref{sec:bg-fpga}: every layer becomes its own hardware block, lowered
through high-level synthesis into a bitstream. The target is the Ultra96, a Zynq
UltraScale+ device, where the programmable logic holds the accelerator and the
processing system streams images to it and reads back predictions.

FINN exposes many build parameters; those that drive the memory and throughput
trade-off we faced are the folding of the compute blocks and their adjacent stream
adapters, controlled mainly through PE and SIMD factors; the depth of each
inter-layer FIFO; the memory style of each block, that is whether its storage
maps to block RAM or to distributed LUT memory, configurable per block and
otherwise left to a default heuristic; and the target clock period. We tune these
in the design-space exploration of Section~\ref{sec:res-dse};
Section~\ref{sec:res-hw} reports the deployed configurations.

Block RAM is worth a closer look, since it turns out to be the resource that bounds
our design. The FPGA fabric used here exposes it as 36-Kb primitives, each usable
as two independent 18-Kb halves and accessed through one write and one read port.
The hardware blocks generated by FINN need storage for weights, thresholds, and
inter-layer streams. Through memory-style attributes, FINN can steer that storage
toward block RAM or distributed LUT memory, depending on the block and backend. In
the design space exploration we exploit this configurability
specifically on FIFOs, where it is a direct way to relieve BRAM pressure without
changing arithmetic. This matters
especially for shallow FIFOs, which can reserve a block-RAM primitive or half
primitive while storing only a small fraction of its capacity. FINN leaves the
FIFO memory style to a default heuristic that sends FIFOs below a depth threshold
to LUT memory and the rest to block RAM, and that assignment is the lever we act
on in Section~\ref{sec:res-dse}.
